\section{Barrier bounds, guessing numbers and index codes}\label{sec:consequences}

The relative form of Theorem~D gives exact values on triangle-free graphs of maximum degree three: the guessing number
is the fractional matching number, and the broadcast rate of index coding is the number of vertices minus it. We end the
section with a general consequence for relations that become one of our classes after deleting a few vertices.

For a graph $G$ and $S\subseteq V$ let $I(S)=\{v\notin S:N_G(v)\subseteq S\}$, the set of isolated vertices of $G-S$
(\lean{isolatedOut}).

\begin{theorem}[Barrier bound]\label{thm:barrier}
\formal{Barrier.lean}{isolatedOut, barrier_ineq, log_card_le_barrier}
Let $G$ be a triangle-free graph of maximum degree at most~$3$ and $S\subseteq V$. Every polymatroid $h$ valid for $N_G$
satisfies
\[
  2h(V)\;\le\;2\sum_{v\in S}h(v)+\sum_{v\notin S\cup I(S)}h(v),
\]
and every storage code $C$ for $G$ over a nonempty finite alphabet $\alpha$ satisfies
$2\log\card C\le(\card V+\card S-\card{I(S)})\log\card\alpha$.
\end{theorem}

\begin{outline}
Let $U=S\cup I(S)$. Theorem~N3 gives a decomposition from the context $U$ (\lean{n3_decomp}), and the relative form of
Theorem~D gives $2(h(V)-h(U))\le\sum_{v\notin U}(h(U\cup\{v\})-h(U))\le\sum_{v\notin U}h(v)$. The vertices of $I(S)$
have all their neighbours in $S$, so $h(U)=h(S)\le\sum_{v\in S}h(v)$ by the closure lemma and subadditivity. For codes,
apply this to the entropy of a uniformly random codeword, as in \cref{thm:codes}.
\end{outline}

\begin{definition}[Fractional matchings]\label{def:fracmatching}
\formal{Matching.lean}{IsFracMatching, fracWeight, fracMatchingNumber}
A \emph{fractional matching} of a graph $G$ is a symmetric function $y\colon V\times V\to\R_{\ge0}$ that vanishes on
non-adjacent pairs and satisfies $\sum_uy(v,u)\le1$ for every $v$. Its \emph{weight} is $\frac12\sum_{u,v}y(u,v)$, each
edge counted once. The \emph{fractional matching number} $\nustar(G)$ is the supremum of the weights.
\end{definition}

\begin{proposition}[The fractional Tutte--Berge formula]\label{prop:fracTB}
\formal{Matching.lean}{fracWeight_le_barrier, exists_doubleCover_matching, dcFrac, dcFrac_isFracMatching, fracWeight_dcFrac, fracMatchingNumber_spec, fracMatchingNumber_eq_min_barrier}
For every finite simple graph $G$, the supremum $\nustar(G)$ is attained by a half-integral fractional matching, and
\[
  2\nustar(G)\;=\;\min_{S\subseteq V}\bigl(\card V+\card S-\card{I(S)}\bigr).
\]
\end{proposition}

The formula is classical (see~\cite{ScheinermanUllman1997}); we formalize it because the next theorem needs both
directions.

\begin{outline}
Weak duality (\lean{fracWeight_le_barrier}): the vertices of $I(S)$ send all their weight into $S$, which receives at most
$\card S$, and every other vertex carries weight at most one, so twice the weight is at most
$\card S+(\card V-\card{I(S)})$. For the converse, take a set $X_0$ of largest deficiency $d=\card{X_0}-\card{N(X_0)}$,
where $N(X)$ is the set of neighbours of $X$. Hall's theorem~\cite{Hall1935}, in its Mathlib form, applied to the
bipartite double cover of $G$ with $d$ dummy vertices added on one side, gives a set $D\subseteq V$ and an injective map
$g$ on $D$ sending each vertex to a neighbour, with at most $d$ vertices outside $D$
(\lean{exists_doubleCover_matching}). Putting weight $\frac12$ on each pair $\{v,g(v)\}$, once per direction, gives a
half-integral fractional matching of weight $\card D/2$ (\lean{dcFrac}, \lean{fracWeight_dcFrac}). For
$S=N(X_0)\setminus X_0$ the set $I(S)$ contains $X_0\setminus N(X_0)$, which gives
$\card D\ge\card V+\card S-\card{I(S)}$. With weak duality, both bounds are equalities (\lean{fracMatchingNumber_spec}).
\end{outline}

\begin{theorem}[The guessing number is the fractional matching number]\label{thm:gn}
\formal{Matching.lean}{log_card_le_fracMatchingNumber, exists_code_fracMatchingNumber, guessingNumber_eq_fracMatching}
\begin{enumerate}[label=(\alph*)]
\item\label{it:gn-lower} For every finite simple graph $G$ and every nonempty finite set $\beta$, some storage code for
$G$ over the alphabet $\beta\times\beta$ has at least $\card{\beta\times\beta}^{\nustar(G)}$ words.
\item\label{it:gn-upper} If $G$ is triangle-free of maximum degree at most~$3$, every storage code for $G$ over every
nonempty finite alphabet $\alpha$ has $\card C\le\card\alpha^{\nustar(G)}$, and over every alphabet $\beta\times\beta$
some storage code has exactly $\card{\beta\times\beta}^{\nustar(G)}$ words.
\end{enumerate}
\end{theorem}

The Lean statements are logarithmic, for instance $\nustar(G)\log\card{\beta\times\beta}\le\log\card C$, with the
convention $\log0=0$. They give the counts above whenever $\nustar(G)>0$ and $\card\beta\ge2$; in the remaining cases the
claims about $\beta\times\beta$ ask only for one word, and a code with a single word is a storage code.

\begin{outline}
\ref{it:gn-lower} Take $D$ and $g$ as in \cref{prop:fracTB}. The code consists of the words $x$ in which, for every
$v\in D$, the second half of $x_v$ equals the first half of $x_{g(v)}$, and all other halves are a fixed symbol: each
vertex of $D$ shares one half of its symbol with its neighbour $g(v)$, and each vertex of $g(D)$ receives one half from
its preimage. Every symbol is then determined by the neighbours, and the $\card\beta^{\card D}=
\card{\beta\times\beta}^{\nustar}$ choices of the shared halves give distinct words
(\lean{exists_code_fracMatchingNumber}).
\ref{it:gn-upper} \Cref{thm:barrier} with a set $S$ attaining the minimum of \cref{prop:fracTB}
(\lean{log_card_le_fracMatchingNumber}); equality by~\ref{it:gn-lower}.
\end{outline}

\begin{corollary}[Guessing numbers]\label{cor:gn}
\formal{Rates.lean}{guessingNumberAt, guessingNumber, isStorageCode_map_equiv, isStorageCode_map_injective, logb_card_le_fracMatchingNumber, guessingNumberAt_le, guessingNumberAt_sq, guessingNumberAt_ge, guessingNumber_eq, tendsto_guessingNumberAt}
For a finite simple graph $G$ and $q\ge2$, let $\gn(G,q)$ be the supremum of $\log_q\card C$ over the storage codes $C$
for $G$ over the alphabet $\{0,\dots,q-1\}$, and let $\gn(G)=\sup_{q\ge2}\gn(G,q)$. If $G$ is triangle-free of maximum
degree at most~$3$, then
\begin{enumerate}[label=(\alph*)]
\item\label{it:gnq} $\gn(G,q)\le\nustar(G)$ for every $q\ge2$;
\item\label{it:gnsq} $\gn(G,r^2)=\nustar(G)$ for every $r\ge2$;
\item\label{it:gnsup} $\gn(G)=\nustar(G)$;
\item\label{it:gnlim} $\gn(G,q)\to\nustar(G)$ as $q\to\infty$.
\end{enumerate}
\end{corollary}

\begin{outline}
The Lean definitions (\lean{guessingNumberAt}, \lean{guessingNumber}) take the supremum with the convention
$\log0=0$ for the empty code, which does not change it, since one-word codes give $0$. \ref{it:gnq} Divide the
logarithmic bound of \cref{thm:gn}\ref{it:gn-upper} by $\log q>0$ (\lean{logb_card_le_fracMatchingNumber},
\lean{guessingNumberAt_le}). \ref{it:gnsq} For $\card\beta=r$, the code of \cref{thm:gn}\ref{it:gn-lower} over
$\beta\times\beta$ has at least $(r^2)^{\nustar(G)}$ words, and a bijection $\beta\times\beta\to\{0,\dots,r^2-1\}$
carries it to a storage code over that alphabet (\lean{isStorageCode_map_equiv}); with~\ref{it:gnq},
$\gn(G,r^2)=\nustar(G)$ (\lean{guessingNumberAt_sq}). \ref{it:gnsup} Every $\gn(G,q)$ is at most $\nustar(G)$, and
$\gn(G,4)=\nustar(G)$ (\lean{guessingNumber_eq}). \ref{it:gnlim} For $q\ge4$ let $r=\lfloor\sqrt q\rfloor\ge2$. An
injection $\{0,\dots,r-1\}^2\to\{0,\dots,q-1\}$ carries the code of \cref{thm:gn}\ref{it:gn-lower} over
$\{0,\dots,r-1\}^2$ to a storage code over $\{0,\dots,q-1\}$ (\lean{isStorageCode_map_injective}), so
$\gn(G,q)\ge\nustar(G)\log(r^2)/\log q$ (\lean{guessingNumberAt_ge}). Since $q<(r+1)^2\le4r^2$, this is at least
$\nustar(G)(1-\log4/\log q)$, using $\nustar(G)\ge0$ (\lean{fracMatchingNumber_nonneg}); with the upper bound
of~\ref{it:gnq}, $\gn(G,q)\to\nustar(G)$ (\lean{tendsto_guessingNumberAt}).
\end{outline}

The restriction to squares in~\ref{it:gnsq} cannot be dropped: by a theorem of Atkins, Rombach and
Skerman~\cite{AtkinsRombachSkerman2017}, for every $q\ge2$ that is not a perfect square and every sufficiently large odd
$n$, the odd cycle $C_n$, which lies in this class and has $\nustar(C_n)=n/2$, has $\gn(C_n,q)=(n-1)/2+\log_qa<n/2$, where
$a$ is the largest divisor of $q$ that is at most $\sqrt q$.

The code of part~\ref{it:gn-lower} realizes Christofides and Markström's fractional clique cover
strategy~\cite{ChristofidesMarkstrom2011}, restated in storage language by Mazumdar, McGregor and
Vorotnikova~\cite{MMV2019}: on a triangle-free graph the cliques are the vertices and the edges, and the fractional clique
cover number is $\card V-\nustar(G)$ (\cref{cor:broadcast}\ref{it:bc-cover}). Christofides and Markström conjectured
that this strategy is optimal on every undirected graph (see~\cite[Conjecture~1]{Dang2015}) and proved it for perfect
graphs and odd cycles~\cite{ChristofidesMarkstrom2011}; Mazumdar, McGregor and Vorotnikova proved it for prisms and
certain chorded cycles~\cite{MMV2019}. The conjecture is false in general: Baber, Christofides, Dang, Vaughan and Riis disproved
it~\cite{Baber2017} (Dang's thesis gives a counterexample on ten vertices~\cite[Section~4.1]{Dang2015}), and Cameron, Dang
and Riis~\cite{CameronDangRiis2016} found triangle-free counterexamples, such as the Higman--Sims graph. For the Clebsch
graph, of degree five, they proved $10\le\gn\le11$, and the fourth paper determined $\gn=10>8=\nustar$
\cite[Corollary~6.11]{FourthPaper2026}. \Cref{thm:gn,cor:gn} prove the conjecture for triangle-free graphs of maximum
degree three, whether the asymptotic guessing number is taken as the supremum over $q$, as here, or as the limit as
$q\to\infty$, in which form the conjecture is often stated (\cref{cor:gn}\ref{it:gnlim}).

\begin{definition}[Index codes]\label{def:indexcode}
\formal{IndexCoding.lean}{IndexCode}
An \emph{index code} for $N$ over an alphabet $\alpha$ with a finite message set $M$ is an encoding map $E\colon\alpha^V\to M$
together with decoders $D_v$, $v\in V$, such that $D_v(E(x),x)=x_v$ for every word $x$, where $D_v(m,x)$ depends only on
$m$ and on the restriction of $x$ to $N(v)$.
\end{definition}

\begin{theorem}[Index codes and fractional matchings]\label{thm:index}
\formal{IndexCoding.lean}{IndexCode.isStorageCode_fibre, log_card_message_ge, exists_indexCode}
\begin{enumerate}[label=(\alph*)]
\item If $G$ is triangle-free of maximum degree at most~$3$, every index code for $N_G$ over a nonempty finite alphabet
$\alpha$ has $\log\card M\ge(\card V-\nustar(G))\log\card\alpha$.
\item For every finite simple graph $G$ and every finite abelian group $\beta$, there is an index code for $N_G$ over
$\beta\times\beta$ with $\log\card M=(\card V-\nustar(G))\log\card{\beta\times\beta}$.
\end{enumerate}
\end{theorem}

\begin{outline}
(a) The words with a common message form a storage code (\lean{IndexCode.isStorageCode_fibre}), as in the fourth
paper~\cite[Remark~4.9(b)]{FourthPaper2026}; each has at most $\card\alpha^{\nustar}$ words by \cref{thm:gn}, and the
$\card M$ fibres cover $\alpha^V$. (b) With $D$ and $g$ as in \cref{prop:fracTB}, broadcast, for every $v$, the
difference between the second half of $x_v$ and the first half of $x_{g(v)}$ when $v\in D$, and the second half of $x_v$
otherwise, and broadcast the first halves of the vertices outside $g(D)$. A vertex recovers the second half of its symbol
from its message and the first half of $x_{g(v)}$, and, if $v=g(u)$, the first half from the message of $u$ and the second
half of $x_u$; both $g(v)$ and $u$ are neighbours. The message set has $\card\beta^{2\card V-\card D}$ elements.
\end{outline}

The duality between storage codes and index codes behind part~(a) is that of Mazumdar~\cite{Mazumdar2014}, Shanmugam and
Dimakis~\cite{ShanmugamDimakis2014}, and Arbabjolfaei and Kim~\cite{ArbabjolfaeiKim2015}.

\begin{corollary}[The broadcast rate and fractional clique covers]\label{cor:broadcast}
\formal{Rates.lean}{IndexCode.ofInjective, IndexCode.restrict, broadcastBits, broadcastBits_nonempty, broadcastBits_ge, broadcastBits_le_succ, broadcastBits_even_le, broadcastBits_even, fracMatchingNumber_nonneg, fracMatchingNumber_le_card, tendsto_broadcastBits, IsFracCover, coverWeight, fracCliqueCoverNumber, fracCliqueCoverNumber_eq}
Let $G$ be a finite simple graph. For $t\ge0$ let $\beta_t$ be the least $n$ such that $N_G$ has an index code over the
alphabet $\{0,1\}^t$ with message set $\{0,1\}^n$.
\begin{enumerate}[label=(\alph*)]
\item\label{it:bc-every} For all $t,s\ge0$, $\beta_t\le\beta_{t+1}$ and $\beta_{2s}\le2(\card V-\nustar(G))s$.
\item\label{it:bc-rate} If $G$ is triangle-free of maximum degree at most~$3$, then $(\card V-\nustar(G))t\le\beta_t$ for
every $t$, $\beta_{2s}=2(\card V-\nustar(G))s$ for every $s$, and $\beta_t/t\to\card V-\nustar(G)$ as $t\to\infty$.
\item\label{it:bc-cover} Call weights $y_v\ge0$ on the vertices and $y(u,v)=y(v,u)\ge0$ on pairs of vertices, vanishing on
non-adjacent pairs, a \emph{fractional cover by vertices and edges} if $y_v+\sum_uy(v,u)\ge1$ for every $v$, and let its
weight be $\sum_vy_v+\frac12\sum_{u,v}y(u,v)$, each edge counted once. The infimum of these weights is
$\card V-\nustar(G)$.
\end{enumerate}
\end{corollary}

\begin{outline}
\ref{it:bc-every} The identity code shows that $\beta_t$ exists (\lean{broadcastBits_nonempty}). An index code over
$\{0,1\}^{t+1}$ restricts to $\{0,1\}^t$: append a fixed bit to every symbol, and drop it after decoding
(\lean{IndexCode.restrict}, \lean{broadcastBits_le_succ}). \Cref{thm:index}(b) with $\beta=\F_2^s$ gives an index code
over $\beta\times\beta$, which is $\{0,1\}^{2s}$ up to a bijection, whose message set has $2^n$ elements with
$n=2(\card V-\nustar(G))s$; a bijection of message sets (\lean{IndexCode.ofInjective}) gives
\lean{broadcastBits_even_le}. \ref{it:bc-rate} \Cref{thm:index}(a) with $\alpha=\{0,1\}^t$ gives the lower bound
(\lean{broadcastBits_ge}), and with~\ref{it:bc-every} the equality for even $t$ (\lean{broadcastBits_even}). For odd $t$,
$\beta_t\le\beta_{t+1}=(\card V-\nustar(G))(t+1)$, so $\card V-\nustar(G)\le\beta_t/t\le(\card V-\nustar(G))(1+1/t)$
for all $t\ge1$, where $0\le\nustar(G)\le\card V$ (\lean{fracMatchingNumber_nonneg}, \lean{fracMatchingNumber_le_card},
\lean{tendsto_broadcastBits}). \ref{it:bc-cover} In Lean the covers are \lean{IsFracCover}, their weight is
\lean{coverWeight}, the infimum is \lean{fracCliqueCoverNumber}, and the claim is \lean{fracCliqueCoverNumber_eq}. Take $S$ attaining the minimum in
\cref{prop:fracTB}, and put $z_v=1$ on $I(S)$, $z_v=0$ on $S$ and $z_v=\frac12$ elsewhere. An edge at a vertex of $I(S)$
ends in $S$, so $z_u+z_v\le1$ on every edge, and every cover has weight at least
$\sum_vz_v(y_v+\sum_uy(v,u))\ge\sum_vz_v=\frac12(\card V-\card S+\card{I(S)})=\card V-\nustar(G)$. Conversely, an optimal
fractional matching $x$ gives the cover with $y(u,v)=x(u,v)$ and $y_v=1-\sum_ux(v,u)$, of weight $\card V-\nustar(G)$.
\end{outline}

The limit in~\ref{it:bc-rate} is the broadcast rate of Blasiak, Kleinberg and Lubetzky~\cite{BKL2013}, so the broadcast
rate of a triangle-free graph of maximum degree three is $\card V-\nustar(G)$. In a triangle-free graph the cliques are
the vertices and the edges, so the infimum in~\ref{it:bc-cover} is the fractional clique cover number, which therefore
equals $\card V-\nustar(G)$. The bound of \cref{thm:index}(b) is thus the fractional clique cover upper bound on the
broadcast rate~\cite{BKL2013}, realized over $\beta\times\beta$ by a double-cover matching. For some graphs of the class
these values were known: Blasiak, Kleinberg and Lubetzky determined the broadcast rate of cycles and of the cubic
circulants, and computed that of the Petersen graph numerically~\cite{BKL2013}.

\paragraph{Why it matters.} For triangle-free subcubic side information, \cref{thm:barrier,thm:gn,thm:index} and
\cref{cor:gn,cor:broadcast} determine the guessing number and the broadcast rate; the bounds hold for all codes and all
alphabets and are attained over the alphabets $\beta\times\beta$. In the language of the fourth paper's crux
(\cref{sec:construction}), on such side information an index code must broadcast at least $\card V-\nustar(G)\ge\card
V/2$ symbols, and the extremal codes are fractional pairings: there is no room for a coding gain.

\subsection{Relations near a class}\label{sec:paid}

All our classes are defined by exact conditions. A few exceptional vertices cost only what they could carry.

\begin{corollary}[Paid deletion]\label{cor:paid}
\formal{PairedDistance.lean}{card_sq_le_of_paidDeletion}
Let $N$ be a recovery relation on $V$ and $Z\subseteq V$. If $N-Z$ satisfies \eqref{eq:star}, then every storage code
$C$ for $N$ over a nonempty finite alphabet $\alpha$ satisfies $\card C^2\le\card\alpha^{\card V+\card Z}$.
\end{corollary}

\begin{outline}
In the formal statement, $N-Z$ is the relation \lean{DirectedRecovery.inducedNeighborhood}~$N\,(V\setminus Z)$ induced
on $V\setminus Z$, as in \emph{Directed Recovery Structures}. By \cref{thm:codes}
(\lean{MotionDepthCodes.card_sq_le_of_rateBound_half}), every storage code for $N-Z$ has at most
$\card\alpha^{(\card V-\card Z)/2}$ words. The paid-conditioning lemma of
\emph{Directed Recovery Structures}~\cite[Lemma~2.2]{DirectedRecovery2026}
(\lean{DirectedRecovery.storageCode_conditioning_bound}) fixes the values on $Z$: every fibre is a storage code for
$N-Z$, and there are at most $\card\alpha^{\card Z}$ fibres, so $\card C\le\card\alpha^{\card Z+(\card V-\card Z)/2}$.
Squaring gives the claim.
\end{outline}

For $\card\alpha\ge2$ the conclusion says that $\log_{\card\alpha}\card C\le(\card V+\card Z)/2$.

Combined with \cref{lem:tt-paired}, the corollary shows that relations with
few short paired paths are as good as TT-free ones: if deleting the loops, the middle vertices of transitive triangles and
a further set $Z'$ leaves a relation in one of our classes, then storage codes have rate at most
$\frac12+(b_E(3)+\card{Z'})/(2\card V)$.

\paragraph{Why it matters.} This is how the exact theorems of this paper apply to the asymptotic setting of the
paired-distance question~\cite[Question~9.1]{DirectedRecovery2026}: a construction must keep its cone relations away
from every class here by a linear number of vertices (\cref{sec:construction}).
